\clearpage\section{Requirements and actors}\label{sec:requirements}
\subsection{Actors and trust boundaries}
An \textbf{author/lecturer} owns application courses and their documents and modules. A local account can draft and export without a confirmed Moodle identity, subject to generation limits. A \textbf{Moodle-authorized lecturer} is the same person with a confirmed identity and authoring capabilities in a specific Moodle course; linking an identity does not grant course permissions.

An \textbf{application administrator} manages accounts and failed jobs and prepares studies. The administrator role is not a general right to inspect another author's documents. An \textbf{invited rater} sees blinded study snapshots and their own ratings. A user may have more than one role, but authorization is checked for the requested action. \textbf{Moodle} is an external publishing and confirmation actor. The \textbf{LLM service} supplies generation and transcription. \textbf{Learners} are stakeholders who use Moodle after publication; they are outside the authoring application's actor boundary.
\subsection{Functional requirements}
\par
Table~\ref{tab:requirements} presents functional requirements.
\begin{longtable}{|p{1.5cm}|p{13.5cm}|}
\caption{Functional requirements.}\label{tab:requirements}\\\hline
ID & Required behavior\\\hline\endfirsthead
\hline ID & Required behavior\\\hline\endhead
FR01 & Register, sign in, renew a session and sign out; disabled users lose access.\\\hline
FR02 & Create and manage owned courses; bind an authorized Moodle target course.\\\hline
FR03 & Upload PDF/DOCX/PPTX, HTML/TXT, SRT/VTT and the audio/video formats listed in Section~\ref{sec:scope} and inspect asynchronous processing states and failures.\\\hline
FR04 & Generate a seven-part module with model, output language and guidance settings.\\\hline
FR05 & Inspect source passages next to generated claims through inline markers.\\\hline
FR06 & Edit lesson components or questions manually and through chat while preserving the original AI version.\\\hline
FR07 & Download generated slides as PPTX.\\\hline
FR08 & Publish an owned RAG module as Moodle Book, Practice and Quiz after capability checks.\\\hline
FR09 & Request and confirm Moodle identity linking without disclosing whether an email exists.\\\hline
FR10 & Inspect and retry failed jobs within role scope.\\\hline
FR11 & Prepare a study, snapshot paired modules and invite active raters.\\\hline
FR12 & Show blinded packets in a stable per-rater shuffled order and save rubric ratings.\\\hline
FR13 & Administer local accounts, including active/admin state and Moodle-account association.\\\hline
FR14 & Ask source-grounded course questions, inspect citations and retrieve conversation history within the owned course.\\\hline
FR15 & List available generation models and save or clear the author's default model preference, validating availability.\\\hline\end{longtable}
\subsection{Non-functional requirements}
The following are design acceptance targets for Phase 2, rather than measured service guarantees. They make isolation, workload and recovery verifiable.
\par
Table~\ref{tab:nfr} presents non-functional requirements and acceptance criteria.
\begin{longtable}{|p{1.6cm}|p{4.0cm}|p{8.3cm}|}\caption{Non-functional requirements and acceptance criteria.}\label{tab:nfr}\\\hline
ID & Concern & Verification criterion\\\hline\endfirsthead
\hline ID & Concern & Verification criterion\\\hline\endhead
NFR01 & Isolation & All cross-owner API and runtime-role database probes are denied; no protected content appears in an unauthorized response.\\\hline
NFR02 & Responsiveness & Under the declared single-worker workload, job acknowledgement p95 is at most 2 s, excluding upload transfer; the client polls every 2--5 s and stops on terminal states.\\\hline
NFR03 & Recoverability & At most three automatic attempts per transient job failure and one corrective retrieval retry; one worker claims a job at a time; retry history remains visible.\\\hline
NFR04 & Source integrity & Every accepted citation resolves to an authorized physical chunk; first-edit tests preserve the original across all later edits.\\\hline
NFR05 & Reproducibility & All formal packets have source/prompt/model/settings hashes; cross-model fallback is disabled; all failures and attempts are retained.\\\hline
NFR06 & Accessibility & Main task flows satisfy Web Content Accessibility Guidelines (WCAG) 2.2 AA, including keyboard completion, visible focus, labelled controls and non-colour-only status; control contrast is at least 3:1.\\\hline
NFR07 & Bounded workload & One heavy background job runs per worker in the baseline configuration. Streamed uploads enforce 50 MB documents, 200 MB audio and 500 MB video; excess inputs return a defined error.\\\hline
NFR08 & Cost observability & Every provider attempt records tokens, served model, elapsed time and status; reconciled totals include failed/retried calls and report local memory separately.\\\hline
\end{longtable}
\needspace{.60\textheight}
\subsection{Permission matrix}

Table~\ref{tab:permissions} separates operational administration from invited rating. An administrator who also authors content retains the same ownership checks as an ordinary author.
\begin{table}[H]\centering\small\caption{Authorization by role; ownership remains required.}\label{tab:permissions}
\begin{tabularx}{\linewidth}{|p{3.4cm}|Y|Y|Y|Y|}\hline
Action & Local author & Linked author & Administrator & Invited rater\\\hline
Draft/edit/export & Own content & Own content & Own content only & Only if separately an author\\\hline
Publish to Moodle & No & Remote course rights & Same remote checks & No rating-role privilege\\\hline
Read other authors' courses & No & No & No general override & No\\\hline
Retry failed jobs & Own & Own & Narrow operational recovery & No\\\hline
Manage accounts / studies & No & No & Yes & No\\\hline
Read/rate blinded packets & If invited & If invited & If invited & Assigned packets; own rating\\\hline
\end{tabularx}\end{table}
\clearpage
\subsection{Requirement traceability}\label{sec:traceability}
Table~\ref{tab:traceability} connects each functional requirement to its use case, related interface figure and API group. The course workspace supplies the course Q\&A entry; the model-selection dialog and account menu supply preference entry points. These associations identify interaction surfaces without treating an authoring chat-edit image as a separate course-answering screen. Detailed endpoint contracts appear in Section~\ref{sec:api}.
\begingroup\small
\begin{longtable}{|p{1.1cm}|p{1.0cm}|p{3.0cm}|p{4.2cm}|p{4.1cm}|}\caption{Functional requirement--use case--screen--API traceability.}\label{tab:traceability}\\\hline
FR & UC & Interface figures & Primary action & API group\\\hline\endfirsthead
\hline FR & UC & Interface figures & Primary action & API group\\\hline\endhead
FR01 & UC01 & \ref{fig:ui-screen-01}, \ref{fig:ui-screen-02} & Account access & \code{auth}\\\hline
FR02 & UC02 & \ref{fig:ui-screen-03}, \ref{fig:ui-screen-04} & Course workspace & \code{courses}\\\hline
FR03 & UC03 & \ref{fig:ui-screen-04}, \ref{fig:ui-screen-14} & Upload/status & \code{documents}\\\hline
FR04 & UC04 & \ref{fig:ui-screen-05} & Generation form & \code{modules/generate}\\\hline
FR05 & UC05 & \ref{fig:ui-screen-06} & Evidence panel & \code{modules}\\\hline
FR06 & UC06 & \ref{fig:ui-screen-06}, \ref{fig:ui-screen-07} & Component/chat edit & \code{modules}\\\hline
FR07 & UC07 & \ref{fig:ui-screen-07} & Slide download & \code{modules}\\\hline
FR08 & UC08 & \ref{fig:ui-screen-08}, \ref{fig:ui-screen-09} & Publication review & \code{modules}\\\hline
FR09 & UC09 & \ref{fig:ui-screen-09}, \ref{fig:ui-screen-15} & Account linking & \code{moodle/account}\\\hline
FR10 & UC11 & \ref{fig:ui-screen-04}, \ref{fig:ui-screen-11} & Failure recovery & \code{jobs; admin/jobs}\\\hline
FR11 & UC12 & \ref{fig:ui-screen-12} & Study setup & \code{studies}\\\hline
FR12 & UC13 & \ref{fig:ui-screen-13} & Blinded rating & \code{studies/mine}\\\hline
FR13 & UC10 & \ref{fig:ui-screen-10} & Account administration & \code{admin/users}\\\hline
FR14 & UC14 & \ref{fig:ui-screen-04} & Course Q\&A entry; source context & \code{courses/{id}/chat}\\\hline
FR15 & UC15 & \ref{fig:ui-screen-05}, \ref{fig:ui-screen-15} & Model selection; account menu & \code{models}\\\hline
\end{longtable}\endgroup
